\ifdefined\gkver\else\def\gkver{student}\fi
\documentclass[\gkver]{gaokaozhenti}
\begin{document}

\chapter{1994年全国}

\begin{enumerate}
\item
%% number: 1
%% paperName: 1994年全国高考第 1 题 3 分
%% typeId: 1
%% type: 单选题
%% chapter: 直线运动
%% point: 运动图象
%% method: 无
%% score: 3
%% degree: 850
%% duplicateId: 0
%% body:
将物体竖直向上抛出后，能正确表示其速率 $v$ 随时间 $t$ 的变化关系的图线是图中图\xzanswer{D}
\fourchoices[answer=D, ispicture=true, h=2.2cm]
{figs/01a.png}
{figs/01b.png}
{figs/01c.png}
{figs/01d.png}

%% answer: D
%% memo:
\memoanswer{
物体做竖直上抛运动，在上升过程中，由 $v_{t}=v_{0}-gt$ 知，上升到最高点时 $v_{t}=0$，然后再做自由落体运动，有 $v_{t}=gt$，故 D 正确。
}

\item
%% number: 2
%% paperName: 1994年全国高考第 2 题 3 分
%% typeId: 1
%% type: 单选题
%% chapter: 气体性质
%% point: 气体图象
%% method: 无
%% score: 3
%% degree: 750
%% duplicateId: 0
%% body:
图中 A、B 两点代表一定质量理想气体的两个不同的状态，状态 A 的温度为 $T_{A}$，状态 B 的温度为 $T_{B}$。由图可知\xzanswer{C}
\onepicture[w=4.5cm]{figs/02.png}
\fourchoices[answer=C]
{$T_{B}=2T_{A}$}
{$T_{B}=4T_{A}$}
{$T_{B}=6T_{A}$}
{$T_{B}=8T_{A}$}

%% answer: C
%% memo:
\memoanswer{
对于一定质量的理想气体，由理想气体状态方程得 $\frac{p_{A}V_{A}}{T_{A}}=\frac{p_{B}V_{B}}{T_{B}}$，解得 $\frac{T_{B}}{T_{A}}=\frac{p_{B}V_{B}}{p_{A}V_{A}}=6$，故 C 正确。
}

\item
%% number: 3
%% paperName: 1994年全国高考第 3 题 3 分
%% typeId: 1
%% type: 单选题
%% chapter: 气体性质
%% point: 物体的内能
%% method: 无
%% score: 3
%% degree: 750
%% duplicateId: 0
%% body:
金属制成的气缸中装有柴油与空气的混合物。有可能使气缸中柴油达到燃点的过程是\xzanswer{A}
\fourchoices[answer=A]
{迅速向里推活塞}
{迅速向外拉活塞}
{缓慢向里推活塞}
{缓慢向外拉活塞}

%% answer: A
%% memo:
\memoanswer{
迅速向气缸内推活塞，气缸内的柴油与空气来不及与气缸和气缸外周围环境进行热交换（可看成绝热压缩），致使温度急剧升高，达到柴油的燃点，故 A 正确。
}

\item
%% number: 4
%% paperName: 1994年全国高考第 4 题 3 分
%% typeId: 1
%% type: 单选题
%% chapter: 曲线运动
%% point: 人造地球卫星
%% method: 无
%% score: 3
%% degree: 850
%% duplicateId: 0
%% body:
人造地球卫星的轨道半径越大，则\xzanswer{B}
\fourchoices[answer=B]
{速度越小，周期越小}
{速度越小，周期越大}
{速度越大，周期越小}
{速度越大，周期越大}

%% answer: B
%% memo:
\memoanswer{
人造地球卫星绕地球做圆周运动的向心力由地球对卫星的万有引力提供，有 $G\frac{Mm}{r^{2}}=m\frac{v^{2}}{r}$，得 $v=\sqrt{\frac{GM}{r}}$，即 $v\propto\sqrt{\frac{1}{r}}$，轨道半径越大，卫星运行的线速度越小。又 $G\frac{Mm}{r^{2}}=m\left(\frac{2\pi}{T}\right)^{2}r$，得 $T=2\pi\sqrt{\frac{r^{3}}{GM}}$，$T\propto\sqrt{r^{3}}$，轨道半径越大，卫星运行的周期越大，最小周期是 $r=R$ 时，$R$ 是地球半径，故 B 正确。
}

\item
%% number: 5
%% paperName: 1994年全国高考第 5 题 3 分
%% typeId: 1
%% type: 单选题
%% chapter: 电路
%% point: 动态分析
%% method: 无
%% score: 3
%% degree: 850
%% duplicateId: 0
%% body:
如图所示的电路中，电源的电动势为 $\varepsilon$，内阻为 $r$。当可变电阻的滑片 P 向 b 点移动时，电压表 V$_{1}$ 的读数 $U_{1}$ 与电压表 V$_{2}$ 的读数 $U_{2}$ 的变化情况是\xzanswer{A}
\onepicture[w=5cm]{figs/05.png}
\fourchoices[answer=A]
{$U_{1}$ 变大，$U_{2}$ 变小}
{$U_{1}$ 变大，$U_{2}$ 变大}
{$U_{1}$ 变小，$U_{2}$ 变小}
{$U_{1}$ 变小，$U_{2}$ 变大}

%% answer: A
%% memo:
\memoanswer{
当可变电阻的滑片 P 向 b 移动时，滑动变阻器电阻变大，电路总电阻变大，干路中电流 $I=\frac{E}{R+r}$ 减小，又路端电压 $U=E-Ir$，故 $U$ 增大，即电压表 V$_{1}$ 的读数 $U_{1}$ 变大，由于电压表 V$_{2}$ 的示数 $U_{2}=IR$，因 $R$ 不变，$I$ 减小，所以 $U_{2}$ 减小，故 A 正确。
}

\item
%% number: 6
%% paperName: 1994年全国高考第 6 题 3 分
%% typeId: 1
%% type: 单选题
%% chapter: 光学
%% point: 光谱
%% method: 无
%% score: 3
%% degree: 850
%% duplicateId: 0
%% body:
太阳的连续光谱中有许多暗线，它们对应着某些元素的特征谱线。产生这些暗线是由于\xzanswer{C}
\fourchoices[answer=C]
{太阳表面大气层中缺少相应的元素}
{太阳内部缺少相应的元素}
{太阳表面大气层中存在着相应的元素}
{太阳内部存在着相应的元素}

%% answer: C
%% memo:
\memoanswer{
太阳表面大气层中存在着某些元素，太阳光照射时就能吸收这些元素光谱对应的光线，太阳的连续光谱中就产生对应着这些元素的特征谱线（暗线），故 C 正确。
}

\item
%% number: 7
%% paperName: 1994年全国高考第 7 题 3 分
%% typeId: 1
%% type: 单选题
%% chapter: 电场
%% point: 电场线
%% method: 无
%% score: 3
%% degree: 650
%% duplicateId: 0
%% body:
若带正电荷的小球只受到电场力作用，则它在任意一段时间内\xzanswer{D}
\fourchoices[answer=D]
{一定沿电力线由高电势处向低电势处运动}
{一定沿电力线由低电势处向高电势处运动}
{不一定沿电力线运动，但一定由高电势处向低电势处运动}
{不一定沿电力线运动，也不一定由高电势处向低电势处运动}

%% answer: D
%% memo:
\memoanswer{
如果带电粒子的初速度不与电力线在一条直线上，则它不一定沿电力线运动，也不一定由高电势处向低电势处运动，故 D 正确。
}

\item
%% number: 8
%% paperName: 1994年全国高考第 8 题 3 分
%% typeId: 1
%% type: 单选题
%% chapter: 电路
%% point: 电功电功率
%% method: 无
%% score: 3
%% degree: 650
%% duplicateId: 0
%% body:
两个电阻，$R_{1}=8\UO$，$R_{2}=2\UO$，并联在电路中。欲使这两个电阻消耗的电功率相等，可行的办法是\xzanswer{A}
\fourchoices[answer=A]
{用一个阻值为 $2\UO$ 的电阻与 $R_{2}$ 串联}
{用一个阻值为 $6\UO$ 的电阻与 $R_{2}$ 串联}
{用一个阻值为 $6\UO$ 的电阻与 $R_{1}$ 串联}
{用一个阻值为 $2\UO$ 的电阻与 $R_{1}$ 串联}

%% answer: A
%% memo:
\memoanswer{
解法一：分析法，电阻消耗的电功率 $P = I^{2}R$，由于电阻 $R_{1}$、$R_{2}$ 并联在电路中，欲使两电阻消耗的电功率相等，$P_{1} = P_{2}$，即 $I_{1}^{2}R_{1} = I_{2}^{2}R_{2}$，得 $I_{1}:I_{2} = \sqrt{R_{2}}:\sqrt{R_{1}} = 1:2$，欲使通过 $R_{2}$ 的电流是 $R_{1}$ 的两倍，$R_{2}$ 支路电阻应等于 $R_{1} = 8\UO$ 电阻的一半，这样需在 $R_{2}$ 支路串联一个 $2\UO$ 的电阻，故 A 正确。

解法二：逐一选项计算，A 项：$R_{\text{串}}=4\UO$，$I_{1}:I_{2}=4:8=1:2$，得 $\frac{P_{1}}{P_{2}}=\left(\frac{I_{1}}{I_{2}}\right)^{2}\cdot\frac{R_{1}}{R_{2}}=1$，故 A 正确；B 项：$R_{\text{串}}=8\UO$，$I_{1}:I_{2}=8:8=1:1$，得 $\frac{P_{1}}{P_{2}}=\left(\frac{I_{1}}{I_{2}}\right)^{2}\cdot\frac{R_{1}}{R_{2}}=4$，故 B 错误；C 项：$R_{\text{串}}=14\UO$，$I_{1}:I_{2}=2:14=\frac{1}{7}$，得 $\frac{P_{1}}{P_{2}}=\left(\frac{I_{1}}{I_{2}}\right)^{2}\cdot\frac{R_{1}}{R_{2}}=\frac{4}{49}$，故 C 错误；D 项：$R_{\text{串}}=10\UO$，$I_{1}:I_{2}=2:10=\frac{1}{5}$，得 $\frac{P_{1}}{P_{2}}=\left(\frac{I_{1}}{I_{2}}\right)^{2}\cdot\frac{R_{1}}{R_{2}}=\frac{4}{25}$，故 D 错误。
}

\item
%% number: 9
%% paperName: 1994年全国高考第 9 题 3 分
%% typeId: 1
%% type: 单选题
%% chapter: 电场
%% point: 电势能电势差电场力做功
%% method: 无
%% score: 3
%% degree: 650
%% duplicateId: 0
%% body:
一个带正电的质点，电量 $q = 2.0 \times 10^{-9}\UC$，在静电场中由 a 点移到 b 点，在这过程中，除电场力外，其他力作的功为 $6.0 \times 10^{-5}\UJ$，质点的动能增加了 $8.0 \times 10^{-5}\UJ$，则 a、b 两点间的电势差 $U_{ab}$ 为\xzanswer{B}
\fourchoices[answer=B]
{$3 \times 10^{4}\UV$}
{$1 \times 10^{4}\UV$}
{$4 \times 10^{4}\UV$}
{$7 \times 10^{4}\UV$}

%% answer: B
%% memo:
\memoanswer{
由动能定理得 $W_{\text{电}} + W_{\text{其他}} = \Delta E_{k}$，所以电场力做功 $W_{\text{电}} = \Delta E_{k} - W_{\text{其他}} = 2.0 \times 10^{-5}\UJ$，又 $W_{\text{电}} = qU$，所以电势差为 $U = \frac{W_{\text{电}}}{q} = 1.0 \times 10^{4}\UV$，故 B 正确。
}

\item
%% number: 10
%% paperName: 1994年全国高考第 10 题 3 分
%% typeId: 1
%% type: 单选题
%% chapter: 磁场
%% point: 带电粒子在磁场中的运动
%% method: 无
%% score: 3
%% degree: 650
%% duplicateId: 0
%% body:
质子和 $\alpha$ 粒子在同一匀强磁场中作半径相同的圆周运动。由此可知质子的动能 $E_{1}$ 和 $\alpha$ 粒子的动能 $E_{2}$ 之比 $E_{1}:E_{2}$ 等于\xzanswer{B}
\fourchoices[answer=B]
{4:1}
{1:1}
{1:2}
{2:1}

%% answer: B
%% memo:
\memoanswer{
由 $r=\frac{mv}{qB}$ 及 $E=\frac{1}{2}mv^{2}$ 得 $E=\frac{1}{2}m\frac{q^{2}B^{2}r^{2}}{m^{2}}=\frac{q^{2}B^{2}r^{2}}{2m}$，所以 $E_{1}:E_{2}=\frac{q_{1}^{2}}{m_{1}}:\frac{q_{2}^{2}}{m_{2}}=1:1$，故 B 正确。
}

\item
%% number: 11
%% paperName: 1994年全国高考第 11 题 3 分
%% typeId: 1
%% type: 单选题
%% chapter: 气体性质
%% point: 气体的状态参量
%% method: 无
%% score: 3
%% degree: 650
%% duplicateId: 0
%% body:
如图所示，一个横截面积为 $S$ 的圆筒形容器竖直放置。金属圆板 A 的上表面是水平的，下表面是倾斜的，下表面与水平面的夹角为 $\theta$，圆板的质量为 $M$。不计圆板与容器内壁之间的摩擦。若大气压强为 $p_{0}$，则被圆板封闭在容器中的气体的压强 $p$ 等于\xzanswer{D}
\onepicture[w=3.5cm]{\includesvg{figs/11}}
\fourchoices[answer=D]
{$p_{0} + \frac{Mg\cos\theta}{S}$}
{$\frac{p_{0}}{\cos\theta} + \frac{Mg}{S\cos\theta}$}
{$p_{0} + \frac{Mg\cos^{2}\theta}{S}$}
{$p_{0} + \frac{Mg}{S}$}

%% answer: D
%% memo:
\memoanswer{
封闭气体对活塞的压力为 $F = \frac{pS}{\cos\theta}$，该力在竖直方向的分力 $F_{\text{竖直}} = F\cos\theta$，活塞受力平衡，有 $p_{0}S + Mg = F_{\text{竖直}}$，解得 $p = p_{0} + \frac{Mg}{S}$，故 D 正确。
}

\item
%% number: 12
%% paperName: 1994年全国高考第 12 题 3 分
%% typeId: 1
%% type: 单选题
%% chapter: 电磁感应
%% point: 电磁感应中图象
%% method: 无
%% score: 3
%% degree: 650
%% duplicateId: 0
%% body:
图中 A 是一边长为 $l$ 的方形线框，电阻为 $R$。今维持线框以恒定的速度 $v$ 沿 $x$ 轴运动，并穿过图中所示的匀强磁场 B 区域。若以 $x$ 轴正方向作为力的正方向，线框在图示位置的时刻作为时间的零点，则磁场对线框的作用力 $F$ 随时间 $t$ 的变化图线为图中的图\xzanswer{B}
\onepicture[w=6cm]{figs/12.png}
\fourchoices[answer=B, ispicture=true, h=2.2cm]
{figs/12a.png}
{figs/12b.png}
{figs/12c.png}
{figs/12d.png}

%% answer: B
%% memo:
\memoanswer{
线圈在第一段位移 $l$ 上，在场区外运动，线框不受安培力作用，$F = 0$；线框经时间 $t = \frac{l}{v}$ 开始进入磁场，并保持 $v$ 不变，线圈右侧切割磁感线，产生电动势 $E = Blv$ 恒定不变，在线圈内产生感应电流 $I$ 恒定不变，线圈受到阻碍相对运动的安培力 $F = BIl$ 大小不变，方向沿 $x$ 轴负方向，取负值；当 $t=\frac{2l}{v}$ 时，线圈全部进入磁场区，回路磁通量不变，感应电流为 0，线框受的安培力为 0，此过程到 $t=\frac{4l}{v}$ 为止；此后线圈以恒定速度出磁场，线圈左侧切割磁感线，产生的电动势 $E=Blv$，感应电流 $I=\frac{E}{R}$ 恒定不变，受到阻碍相对运动的安培力 $F=BIl$ 大小不变，方向仍沿 $x$ 轴负方向，此过程到 $t=\frac{5l}{v}$ 为止，故 B 正确。
}

\item
%% number: 13
%% paperName: 1994年全国高考第 13 题 3 分
%% typeId: 1
%% type: 单选题
%% chapter: 机械振动
%% point: 振动图象
%% method: 无
%% score: 3
%% degree: 450
%% duplicateId: 0
%% body:
图（a）是演示简谐振动图象的装置。当盛沙漏斗下面的薄木板 N 被匀速地拉出时，摆动着的漏斗中漏出的沙在板上形成曲线，该曲线显示出摆的位移随时间变化的关系，板上的直线 OO$_{1}$ 代表时间轴。图（b）表示两个摆中的沙在各自木板上形成的曲线，若板 N$_{1}$ 和板 N$_{2}$ 拉动的速度 $v_{1}$ 和 $v_{2}$ 的关系为 $v_{2} = 2v_{1}$，则板 N$_{1}$、N$_{2}$ 上曲线所代表的振动的周期 $T_{1}$ 和 $T_{2}$ 的关系为\xzanswer{D}
\twopicture[h=3.2cm]
{figs/13a.png}
{figs/13b.png}
\fourchoices[answer=D]
{$T_{2} = T_{1}$}
{$T_{2} = 2T_{1}$}
{$T_{2} = 4T_{1}$}
{$T_{2} = \frac{T_{1}}{4}$}

%% answer: D
%% memo:
\memoanswer{
木板从沙摆下经过，所需时间 $t=\frac{L}{v}$，其中 $L$ 为板长，$v$ 为板的拉动速度，又木板 $N_{1}$ 经过沙摆过程中，沙摆做了一次全振动，木板 $N_{2}$ 经过沙摆过程中，沙摆做了两次全振动。因此有 $t_{1}=\frac{L}{v_{1}}$ 且 $t_{1}=T_{1}$；$t_{2}=\frac{L}{v_{2}}$ 且 $t_{2}=2T_{2}$，所以 $\frac{T_{1}}{T_{2}}=\frac{t_{1}}{\frac{1}{2}t_{2}}=\frac{2v_{2}}{v_{1}}=4$，即 $T_{2}=\frac{1}{4}T_{1}$，故 D 正确。
}

\item
%% number: 14
%% paperName: 1994年全国高考第 14 题 5 分
%% typeId: 2
%% type: 多选题
%% chapter: 电路
%% point: 电容
%% method: 无
%% score: 5
%% degree: 650
%% duplicateId: 0
%% body:
连接在电池两极上的平行板电容器，当两极板间的距离减小时，则\xzanswer{ABD}
\fourchoices[answer=ABD]
{电容器的电容 $C$ 变大}
{电容器极板的带电量 $Q$ 变大}
{电容器两极板间的电势差 $U$ 变大}
{电容器两极板间的电场强度 $E$ 变大}

%% answer: ABD
%% memo:
\memoanswer{
由 $C=\frac{\varepsilon S}{4\pi kd}$ 知，当两极板间的距离 $d$ 减小时，$C$ 增大，由于两极板连接在电池两极上，所以两极板间电势差 $U$ 保持不变，电容器所带电量 $Q=UC$ 变大，两极板间的电场强度 $E=\frac{U}{d}$ 变大，故 ABD 正确。
}

\item
%% number: 15
%% paperName: 1994年全国高考第 15 题 5 分
%% typeId: 1
%% type: 单选题
%% chapter: 运动和力
%% point: 动量和冲量
%% method: 无
%% score: 5
%% degree: 850
%% duplicateId: 0
%% body:
若物体在运动过程中受到的合外力不为零，则\xzanswer{B}
\fourchoices[answer=B]
{物体的动能不可能总是不变的}
{物体的动量不可能总是不变的}
{物体的加速度一定变化}
{物体的速度方向一定变化}

%% answer: B
%% memo:
\memoanswer{
例如匀速圆周运动，物体在运动过程中受到的合外力不为零，但物体的动能总是不变的，故 A 错误；由动量定理得 $\Delta p = F \Delta t$，只要 $F$ 不为 0，$\Delta p$ 就不可能总为 0，故 B 正确；只要 $F$ 为恒量，则加速度就不变化，故 C 错误；例如在自由落体运动中，$F$ 不为 0，但速度方向不变化，故 D 错误。
}

\item
%% number: 16
%% paperName: 1994年全国高考第 16 题 3 分
%% typeId: 2
%% type: 多选题
%% chapter: 运动和力
%% point: 连接体
%% method: 无
%% score: 3
%% degree: 550
%% duplicateId: 0
%% body:
如图所示，C 是水平地面，A、B 是两个长方形物块，$F$ 是作用在物块 B 上沿水平方向的力，物体 A 和 B 以相同的速度作匀速直线运动。由此可知，A、B 间的滑动摩擦系数 $\mu_{1}$ 和 B、C 间的滑动摩擦系数 $\mu_{2}$ 有可能是\xzanswer{BD}
\onepicture[w=4.5cm]{\includesvg{figs/16}}
\fourchoices[answer=BD]
{$\mu_{1} = 0$，$\mu_{2} = 0$}
{$\mu_{1} = 0$，$\mu_{2} \neq 0$}
{$\mu_{1} \neq 0$，$\mu_{2} = 0$}
{$\mu_{1} \neq 0$，$\mu_{2} \neq 0$}

%% answer: BD
%% memo:
\memoanswer{
以 A 和 B 一起作为研究对象，A 和 B 以相同的速度做匀速直线运动，则在水平方向上受力平衡，故 C 对 B 的摩擦力 $F_{f} = F$，即 B 与 C 间动摩擦因数 $\mu_{2} \neq 0$，故 AC 错误；A 与 B 相对静止，无相对运动趋势，其间无摩擦力，不能断定 $\mu_{1}$ 是否为 0，可能 $\mu_{1} = 0$ 或 $\mu_{1} \neq 0$，故 BD 正确。
}

\item
%% number: 17
%% paperName: 1994年全国高考第 17 题 5 分
%% typeId: 2
%% type: 多选题
%% chapter: 交流电
%% point: LC振荡回路
%% method: 无
%% score: 5
%% degree: 550
%% duplicateId: 0
%% body:
如图为 LC 振荡电路中电容器极板上的电量 $q$ 随时间 $t$ 变化的图线，由图可知\xzanswer{ACD}
\onepicture[w=4.5cm]{\includesvg{figs/17}}
\fourchoices[answer=ACD]
{在 $t_{1}$ 时刻，电路中的磁场能最小}
{从 $t_{1}$ 到 $t_{2}$，电路中的电流值不断变小}
{从 $t_{2}$ 到 $t_{3}$，电容器不断充电}
{在 $t_{4}$ 时刻，电容器的电场能最小}

%% answer: ACD
%% memo:
\memoanswer{
由 LC 电路的充、放电规律及 $q-t$ 图象可知，$t_{1}$ 时刻电量最大、电场能最大，电流为零、磁场能最小；从 $t_{1}\sim t_{2}$ 过程为放电过程，电量减少、电场能减小，电流增大、磁场能增大；从 $t_{2}\sim t_{3}$ 为反向充电过程，电量增大、电场能增大，电流减小，磁场能减小；在 $t_{4}$ 时刻电量为零，电场能为零，故 ACD 正确。
}

\item
%% number: 18
%% paperName: 1994年全国高考第 18 题 5 分
%% typeId: 2
%% type: 多选题
%% chapter: 机械波
%% point: 波的图象
%% method: 无
%% score: 5
%% degree: 550
%% duplicateId: 0
%% body:
如图所示，在 $xy$ 平面内有一沿 $x$ 轴正方向传播的简谐横波，波速为 $1\Ums$，振幅为 $4\Ucm$，频率为 $2.5\UHz$。在 $t=0$ 时刻，P 点位于其平衡位置上方最大位移处，则距 P 为 $0.2\Um$ 的 Q 点\xzanswer{BD}
\onepicture[w=4.5cm]{\includesvg{figs/18}}
\fourchoices[answer=BD]
{在 $0.1\Us$ 时的位移是 $4\Ucm$}
{在 $0.1\Us$ 时的速度最大}
{在 $0.1\Us$ 时的速度向下}
{在 0 到 $0.1\Us$ 时间内的路程是 $4\Ucm$}

%% answer: BD
%% memo:
\memoanswer{
由题可知，$\lambda=\frac{v}{f}=0.4\Um$，$T=\frac{1}{f}=0.4\Us$，波向右传播（$x$ 轴正向），振幅 $A=4\Ucm$，$s_{PQ}=0.2\Um=\frac{\lambda}{2}$，即 PQ 两点相距半个波长，$t=0$ 时，当 P 点位于平衡位置上方最大位移处，则 Q 点位于平衡位置下方最大位移处，$t=0.1\Us=\frac{1}{4}T$ 时，波向右传播 $\frac{\lambda}{4}$，

\onepicture[w=7cm]{figs/18_an.jpg}

如图所示，从图可知，Q 点在 $t=0.1\Us$ 时刻，位移为 0，振动速度达到最大值，方向向上，故 B 正确 AC 错误；Q 点在 0 到 $0.1\Us$ 时间内的路程等于振幅，为 $4\Ucm$，故 D 正确。
}

\item
%% number: 19
%% paperName: 1994年全国高考第 19 题 5 分
%% typeId: 2
%% type: 多选题
%% chapter: 电场
%% point: 带电粒子在交变电场中的运动
%% method: 无
%% score: 5
%% degree: 450
%% duplicateId: 0
%% body:
图中 A、B 是一对平行的金属板。在两板间加上一周期为 $T$ 的交变电压 $u$。A 板的电势 $\varphi_{A} = 0$，B 板的电势 $\varphi_{B}$ 随时间的变化规律为：在 0 到 $T/2$ 的时间内，$\varphi_{B} = U_{0}$（正的常数）；在 $T/2$ 到 $T$ 的时间内，$\varphi_{B} = -U_{0}$；在 $T$ 到 $3T/2$ 的时间内，$\varphi_{B} = U_{0}$；在 $3T/2$ 到 $2T$ 的时间内，$\varphi_{B} = -U_{0}\ldots\ldots$，现有一电子从 A 板上的小孔进入两板间的电场区内。设电子的初速度和重力的影响均可忽略，则\xzanswer{AB}
\onepicture[w=3.5cm]{figs/19.png}
\fourchoices[answer=AB]
{若电子是在 $t=0$ 时刻进入的，它将一直向 B 板运动}
{若电子是在 $t = T/8$ 时刻进入的，它可能时而向 B 板运动，时而向 A 板运动，最后打在 B 板上}
{若电子是在 $t = 3T/8$ 时刻进入的，它可能时而向 B 板运动，时而向 A 板运动，最后打在 B 板上}
{若电子是在 $t = T/2$ 时刻进入的，它可能时而向 B 板、时而向 A 板运动}

%% answer: AB
%% memo:
\memoanswer{
（1）若电子是在 $t=0$ 时刻进入的，电子在 $0\sim\frac{T}{2}$ 段时间内，将向 B 板做初速度为零的匀加速直线运动；在 $\frac{T}{2}\sim T$ 段时间内，电子虽受力反向，但仍向 B 板做匀减速直线运动，$T$ 时刻速度为零，然后又重复以上运动，因此在该情况下，电子将一直向 B 板运动，故 A 正确；

（2）若电子是在 $t=\frac{T}{8}$ 时刻进入的，则在 $\frac{T}{8}\sim\frac{T}{2}$ 段时间内电子将一直向 B 板做匀加速直线运动；在 $\frac{T}{2}\sim\frac{7T}{8}$ 段时间内，电子受力反向，一直向 B 板做匀减速直线运动；在 $t=\frac{7T}{8}$ 时刻，速度减为零；电子在 $\frac{7T}{8}\sim T$ 段时间内，将向 A 板做匀加速直线运动；在 $T\sim\frac{9}{8}T$ 段时间内，将向 A 板做匀减速直线运动；在 $t=\frac{9}{8}T$ 时刻，电子速度减为零，电子在这段时间内，电子向 B 板运动的时间和位移比折回 A 板运动的时间和位移大，因此，电子可能时而向 B 板运动，时而向 A 板运动，最后打在 B 板上，故 B 正确；

（3）若电子是在 $t=\frac{3}{8}T$ 时刻进入，则电子先向 B 板运动，然后再折向 A 板运动，但向 A 板运动到小孔时，速度仍在增大，故电子只能从 A 孔返回，不能打在 B 板上，故 C 错误；

（4）若电子是在 $t=\frac{1}{2}T$ 时刻进入电场，电子开始就受力向 A 板返回，不能在电场中运动，故 D 错误。
}

\item
%% number: 20
%% paperName: 1994年全国高考第 20 题 5 分
%% typeId: 3
%% type: 填空题
%% chapter: 运动和力
%% point: 动量和能量
%% method: 无
%% score: 5
%% degree: 650
%% duplicateId: 0
%% body:
质量为 $4.0\Ukg$ 的物体 A 静止在水平桌面上。另一个质量为 $2.0\Ukg$ 的物体 B 以 $5.0\Ums$ 的水平速度与物体 A 相撞，碰撞后物体 B 以 $1.0\Ums$ 的速度反向弹回。相撞过程中损失的机械能是 \tkanswer[1.2cm]{6}$\UJ$。

%% answer: 
\jdanswer{
6
}

%% memo:
\memoanswer{
两物体碰撞时，由动量守恒得 $m_{B}v_{B} = m_{A}v_{A}' + m_{B}v_{B}'$，解得 $v_{A}' = \frac{m_{B}v_{B} - m_{B}v_{B}'}{m_{A}} = 3.0\Ums$，相撞过程中损失的机械能为 $\Delta E_{k} = \frac{1}{2}m_{B}v_{B}^{2} - \left(\frac{1}{2}m_{A}v_{A}'^{2} + \frac{1}{2}m_{B}v_{B}'^{2}\right) = 6\UJ$。
}

\item
%% number: 21
%% paperName: 1994年全国高考第 21 题 5 分
%% typeId: 3
%% type: 填空题
%% chapter: 原子和原子核
%% point: 结合能
%% method: 无
%% score: 5
%% degree: 750
%% duplicateId: 0
%% body:
一个铀核衰变为钍核时释放出一个 $\alpha$ 粒子。已知铀核的质量为 $3.853131 \times 10^{-25}\Ukg$，钍核的质量为 $3.786567 \times 10^{-25}\Ukg$，$\alpha$ 粒子的质量为 $6.64672 \times 10^{-27}\Ukg$。在这个衰变过程中释放出的能量等于 \tkanswer[2.4cm]{$8.7 \times 10^{-13}$}$\UJ$。（保留两位数字）

%% answer: 
\jdanswer{
$8.7 \times 10^{-13}$（答 $9.0 \times 10^{-13}$ 到 $8.6 \times 10^{-13}$ 都算对）
}

%% memo:
\memoanswer{
在铀核衰变过程中，亏损的质量为 $\Delta m = m_{\mathrm{U}} - (m_{\mathrm{Th}} + m_{\alpha}) = 9.7 \times 10^{-30}\Ukg$，由质能方程得 $\Delta E = \Delta m c^{2} \approx 8.7 \times 10^{-13}\UJ$。
}

\item
%% number: 22
%% paperName: 1994年全国高考第 22 题 5 分
%% typeId: 4
%% type: 实验题
%% chapter: 力物体的平衡
%% point: 验证平行四边形实验
%% method: 无
%% score: 5
%% degree: 750
%% duplicateId: 0
%% body:
将橡皮筋的一端固定在 A 点，另一端拴上两根细绳，每根细绳分别连着一个量程为 $5\UN$、最小刻度为 $0.1\UN$ 的弹簧测力计。沿着两个不同的方向拉弹簧测力计。当橡皮筋的活动端拉到 O 点时，两根细绳相互垂直，如图所示。这时弹簧测力计的读数可从图中读出。
\begin{enumerate}
	\item
	由图可读得两个相互垂直的拉力的大小分别为 \tkanswer[1cm]{4.0}$\UN$ 和 \tkanswer[1cm]{2.5}$\UN$。（只须读到 $0.1\UN$）
	\item
	在本题的虚线方格纸上按作图法的要求画出这两个力及它们的合力。
\end{enumerate}
\onepicture[align=right, w=6cm]{figs/22.png}

%% answer: 
\jdanswer{
\begin{enumerate}
	\item
	4.0，2.5

	\item
	如图所示。

	\onepicture[w=6cm]{figs/22_an.jpg}
\end{enumerate}
}

%% memo:
\memoanswer{
读弹簧秤示数时，应注意首先找零刻度，尤其是竖直的那个弹簧秤是倒置，它的读数是 $2.5\UN$（而不是 $3.5\UN$），水平放置的弹簧秤读数是 $4\UN$，作图时先标坐标。
}

\item
%% number: 23
%% paperName: 1994年全国高考第 23 题 5 分
%% typeId: 3
%% type: 填空题
%% chapter: 交流电
%% point: 交流电
%% method: 无
%% score: 5
%% degree: 550
%% duplicateId: 0
%% body:
如图所示，一个面积为 $S$ 的矩形线圈在匀强磁场中以其一条边为转轴作匀速转动，磁场方向与转轴垂直，线圈中感应电动势 $e$ 与时间 $t$ 的关系如图所示。感应电动势最大值和周期可由图中读出。则磁感应强度 $B =$ \tkanswer[1.5cm]{$\frac{E_{m}T}{2\pi S}$}。在 $t = \frac{1}{12}T$ 时刻，线圈平面与磁感应强度的夹角等于 \tkanswer[1.5cm]{$30^{\circ}$}。
\onepicture[align=right, w=5cm]{\includesvg{figs/23}}

%% answer: 
\jdanswer{
$\frac{E_{m}T}{2\pi S}$，$30^{\circ}$（答 $150^{\circ}$ 也算对）
}

%% memo:
\memoanswer{
由 $e-t$ 图知，$t=0$ 时，$e=E_{m}$，此时线圈平面与磁场的夹角为零，有 $E_{m}=B\omega S$，又 $\omega=\frac{2\pi}{T}$，联立解得 $B=\frac{E_{m}}{S\omega}=\frac{E_{m}T}{2\pi S}$。

任意时刻 $t$ 时，有 $e=E_{m}\cos\omega t=E_{m}\cos\frac{2\pi}{T}t.$ 当 $t=\frac{T}{12}$ 时，有 $e=E_{m}\cos\left(\frac{2\pi}{T}\times\frac{T}{12}\right)=E_{m}\cos\frac{\pi}{6}$，即 $t=\frac{T}{12}$ 时，线圈转过的角度 $\theta=\omega t$，即 $30^{\circ}$。
}

\item
%% number: 24
%% paperName: 1994年全国高考第 24 题 6 分
%% typeId: 4
%% type: 实验题
%% chapter: 光学
%% point: 测定玻璃折射率
%% method: 无
%% score: 6
%% degree: 650
%% duplicateId: 0
%% body:
如图所示，用三棱镜做测定玻璃折射率的实验。先在白纸上放好三棱镜，在棱镜的一侧插上两枚大头针 P$_{1}$ 和 P$_{2}$，然后在棱镜的另一侧观察，调整视线使 P$_{1}$ 的像被 P$_{2}$ 挡住。接着在眼睛所在的一侧插两枚大头针 P$_{3}$、P$_{4}$，使 P$_{3}$ 挡住 P$_{1}$、P$_{2}$ 的像。P$_{4}$ 挡住 P$_{3}$ 和 P$_{1}$、P$_{2}$ 的像，在纸上标出的大头针位置和三棱镜轮廓如图所示。
\begin{enumerate}
	\item
	在本题的图上画出所需的光路。
	\item
	为了测出棱镜玻璃的折射率，需要测量的量是 \tkanswer[1.5cm]{入射角 $i$}，\tkanswer[1.5cm]{折射角 $r$}，在图上标出它们。
	\item
	计算折射率的公式是 $n =$ \tkanswer[2cm]{$\frac{\sin i}{\sin r}$}。
\end{enumerate}
\onepicture[align=right, w=5cm]{figs/24.png}

%% answer: 
\jdanswer{
\begin{enumerate}
	\item
	如图所示，画出通过 P$_{1}$、P$_{2}$、P$_{3}$、P$_{4}$ 的光线，包括在棱镜内的那部分路径。

	\onepicture[w=7cm]{figs/24_an.jpg}

	\item
	入射角 $i$ 和折射角 $r$（见图）（或线段 EF、OE、GH、OG）

	\item
	$n = \frac{\sin i}{\sin r}$（或 $n = \frac{EF}{GH}\cdot\frac{OG}{OE}$）
\end{enumerate}
}

%% memo:
\memoanswer{
根据插针确定光线路径画出光路图。连接 $P_{1}$、$P_{2}$ 点作直线交 AC 边于 O 点，连接 $P_{3}$、$P_{4}$ 点作直线交 AB 边于 $O'$ 点，再连接 O、$O'$ 点，即光线传播路径。过 O 点作 AC 的法线 FH，可量出 $i$、$r$，由折射定律得 $n=\frac{\sin i}{\sin r}$；用量角器量角度误差较大或不方便，可以转化成量长度，因 $\sin i = \frac{EF}{EO}$，$\sin r = \frac{GH}{GO}$，故 $n = \frac{EF \cdot GO}{EO \cdot GH}$；若作图时使 $GO = EO$，则 $n = \frac{EF}{GH}$，可以简化计算。
}

\item
%% number: 25
%% paperName: 1994年全国高考第 25 题 6 分
%% typeId: 4
%% type: 实验题
%% chapter: 电路
%% point: 伏安法测电阻
%% method: 无
%% score: 6
%% degree: 650
%% duplicateId: 0
%% body:
图为用伏安法测量一个定值电阻阻值的实验所需的器材实物图，器材规格如下：
\begin{enumerate}
	\item
	待测电阻 $R_{x}$（约 $100\UO$）；
	\item
	直流毫安表（量程 $0 \sim 10\UmA$，内阻 $50\UO$）；
	\item
	直流电压表（量程 $0 \sim 3\UV$，内阻 $5\UkO$）；
	\item
	直流电源（输出电压 $4\UV$，内阻可不计）；
	\item
	滑动变阻器（阻值范围 $0 \sim 15\UO$，允许最大电流 $1\UA$）；
	\item
	电键一个，导线若干条。
\end{enumerate}
根据器材的规格和实验要求，在本题的实物图上连线。
\onepicture[align=right, w=7cm]{figs/25.png}

%% answer: 
\jdanswer{
如图所示。

\onepicture[w=8cm]{figs/25_an.jpg}
}

%% memo:
\memoanswer{
待测电阻 $R_{x}$ 与电流表内阻 $R_{A}$、电压表内阻 $R_{V}$ 之比分别为 $R_{x}:R_{A}=100:50=2:1$，$R_{x}:R_{V}=100:5\times10^{3}=1:50$，由此可知，直流毫安表应采取“外接法”，电路中最大电流为 $I_{m}=\frac{U}{R_{x}+R_{A}}=\frac{4}{100+50}\UA \approx 27\UmA$，大于直流毫安表的量程 $10\UmA$，故应将滑动变阻器接成分压电路，通过滑动变阻器的最大电流 $I=\frac{U}{R}=\frac{4}{15}\UA=0.27\UA$，小于滑动变阻器允许的最大电流 $1\UA$，滑动变阻器能正常工作。
}

\item
%% number: 26
%% paperName: 1994年全国高考第 26 题 5 分
%% typeId: 4
%% type: 实验题
%% chapter: 直线运动
%% point: 游标卡尺和螺旋测微仪
%% method: 无
%% score: 5
%% degree: 850
%% duplicateId: 0
%% body:
游标卡尺的主尺最小分度为 $1\Umm$，游标上有 20 个小的等分刻度。用它测量一工件的内径，如图所示。该工件的内径为 \tkanswer[1.5cm]{23.85}$\Umm$。
\onepicture[align=right, w=7cm]{figs/26.png}

%% answer: 
\jdanswer{
23.85
}

%% memo:
\memoanswer{
主尺读出 $23\Umm$，游标每一分度是 $0.05\Umm$，第 17 条刻度线与主尺上毫米刻度线对齐，所以游标读数是 $17 \times 0.05\Umm = 0.85\Umm$，故读数为 $23\Umm + 17 \times 0.05\Umm = 23.85\Umm$。
}

\item
%% number: 27
%% paperName: 1994年全国高考第 27 题 5 分
%% typeId: 3
%% type: 填空题
%% chapter: 运动和力
%% point: 动量定理
%% method: 无
%% score: 5
%% degree: 650
%% duplicateId: 0
%% body:
一质量为 $100\Ug$ 的小球从 $0.80\Um$ 高处自由下落到一厚软垫上。若从小球接触软垫到小球陷至最低点经历了 $0.20\Us$，则这段时间内软垫对小球的冲量为 \tkanswer[2cm]{$0.6\UNs$}。（取 $g = 10\Umsq$，不计空气阻力）

%% answer: 
\jdanswer{
$0.6\UNs$
}

%% memo:
\memoanswer{
小球从刚着地到静止，以竖直向上方向为正方向，对这一过程应用动量定理得 $Nt - mgt = 0 - (-mv_{1})$，其中 $v_{1}$ 为小球刚着地时的速度 $v_{1} = \sqrt{2gh} = 4\Ums$，故软垫对小球的冲量为 $Nt = mgt + mv_{1} = 0.6\UNs$。
}

\item
%% number: 28
%% paperName: 1994年全国高考第 28 题 7 分
%% typeId: 6
%% type: 计算题
%% chapter: 光学
%% point: 透镜
%% method: 无
%% score: 7
%% degree: 650
%% duplicateId: 0
%% body:
蜡烛距光屏 $90\Ucm$，要使光屏上呈现出放大到 2 倍的蜡烛像，应选焦距是多大的凸透镜？

%% answer: 
\jdanswer{
$20\Ucm$
}

%% memo:
\memoanswer{
由放大率公式和本题要求，有 $v=2u$，以题给条件 $u + v = 90\Ucm$，解得 $u = 30\Ucm$，$v = 60\Ucm$，由成像公式得 $\frac{1}{u}+\frac{1}{v}=\frac{1}{f}$，解得凸透镜的焦距 $f = 20\Ucm$。
}

\item
%% number: 29
%% paperName: 1994年全国高考第 29 题 10 分
%% typeId: 6
%% type: 计算题
%% chapter: 气体性质
%% point: 玻意耳定律
%% method: 无
%% score: 10
%% degree: 450
%% duplicateId: 0
%% body:
如图所示，可沿气缸壁自由活动的活塞将密封的圆筒形气缸分隔成 A、B 两部分。活塞与气缸顶部有一弹簧相连。当活塞位于气缸底部时弹簧恰好无形变。开始时 B 内充有一定量的气体，A 内是真空。B 部分高度为 $L_{1} = 0.10\Um$。此时活塞受到的弹簧作用力与重力的大小相等。现将整个装置倒置，达到新的平衡后 B 部分的高度 $L_{2}$ 等于多少？设温度不变。
\onepicture[align=right, w=3.5cm]{figs/29.png}

%% answer: 
\jdanswer{
$L_{2} = 0.2\Um$
}

%% memo:
\memoanswer{
设开始时 B 中压强为 $p_{1}$，气缸倒置达到平衡后 B 中压强为 $p_{2}$。分析活塞受力得 $p_{1}S = kL_{1} + Mg$ ①，$p_{2}S + Mg = kL_{2}$ ②，其中 $S$ 为气缸横截面积，$M$ 为活塞质量，$k$ 为弹簧的倔强系数。由题给条件有 $kL_{1}=Mg$ ③。由玻意耳定律，$p_{1}L_{1}=p_{2}L_{2}$ ④。由以上各式得到 $L_{2}^{2}-L_{1}L_{2}-2L_{1}^{2}=0$ ⑤，解得 $L_{2}=2L_{1}=0.2\Um$ ⑥。
}

\item
%% number: 30
%% paperName: 1994年全国高考第 30 题 10 分
%% typeId: 6
%% type: 计算题
%% chapter: 运动和力
%% point: 牛顿定律的应用
%% method: 无
%% score: 10
%% degree: 550
%% duplicateId: 0
%% body:
如图所示，质量 $M = 10\Ukg$ 的木楔 ABC 静置于粗糙水平地面上，滑动摩擦系数 $\mu = 0.02$。在木楔的倾角 $\theta$ 为 $30^{\circ}$ 的斜面上，有一质量 $m = 1.0\Ukg$ 的物块由静止开始沿斜面下滑。当滑行路程 $s = 1.4\Um$ 时，其速度 $v = 1.4\Ums$。在这过程中木楔没有动。求地面对木楔的摩擦力的大小和方向。（重力加速度取 $g = 10\Umsq$）
\onepicture[align=right, w=5cm]{\includesvg{figs/30}}

%% answer: 
\jdanswer{
$f_{2} = 0.61\UN$，方向与图中所设的一致（由 C 指向 B 的方向）
}

%% memo:
\memoanswer{
由匀加速运动的公式 $v^{2}=v_{0}^{2}+2as$，得物块沿斜面下滑的加速度为 $a=\frac{v^{2}}{2s}=0.7\Umsq$。

由于 $a < g\sin\theta = 5\Umsq$，可知物块受到摩擦力作用。

分析物块受力，它受三个力，如图所示，对于沿斜面的方向和垂直于斜面的方向，由牛顿定律，有

\onepicture[w=5cm]{figs/30_an.jpg}

$mg\sin\theta - f_{1}' = ma$ ①

$mg\cos\theta - N_{1}' = 0$ ②

分析木楔受力，它受五个力作用，如图所示，对于水平方向，由牛顿定律，对于水平方向有

$f_{2} + f_{1}\cos\theta - N_{1}\sin\theta = 0$ ③

联立解得地面作用于木楔的摩擦力

$f_{2}=N_{1}\sin\theta-f_{1}\cos\theta=mg\cos\theta\sin\theta-(mg\sin\theta-ma)\cos\theta=ma\cos\theta=0.7\times\frac{2}{3}\UN=0.61\UN$

此力的方向与图中所设的一致（由 C 指向 B 的方向）
}

\item
%% number: 31
%% paperName: 1994年全国高考第 31 题 12 分
%% typeId: 6
%% type: 计算题
%% chapter: 磁场
%% point: 带电粒子在磁场中的运动
%% method: 无
%% score: 12
%% degree: 350
%% duplicateId: 0
%% body:
如图所示，一带电质点，质量为 $m$，电量为 $q$，以平行于 $Ox$ 轴的速度 $v$ 从 $y$ 轴上的 a 点射入图中第一象限所示的区域。为了使该质点能从 $x$ 轴上的 b 点以垂直于 $Ox$ 轴的速度 $v$ 射出，可在适当的地方加一个垂直于 $xy$ 平面、磁感应强度为 $B$ 的匀强磁场。若此磁场仅分布在一个圆形区域内，试求这圆形磁场区域的最小半径。重力忽略不计。
\onepicture[align=right, w=5cm]{\includesvg{figs/31}}

%% answer: 
\jdanswer{
$r = \frac{\sqrt{2}mv}{2qB}$
}

%% memo:
\memoanswer{
质点在磁场中作半径为 $R$ 的圆周运动，由洛伦兹力提供向心力得 $qvB = m\frac{v^{2}}{R}$，解得 $R = \frac{mv}{qB}$ ①。

根据题意，质点在磁场区域中的轨道是半径等于 $R$ 的圆上的 $\frac{1}{4}$ 圆周，这段圆弧应与入射方向的速度、出射方向的速度相切。过 a 点作平行于 $x$ 轴的直线，过 b 点作平行于 $y$ 轴的直线，则与这两直线均相距 $R$ 的 $O'$ 点就是圆周的圆心。质点在磁场区域中的轨道就是以 $O'$ 为圆心、$R$ 为半径的圆（图中虚线圆）上的圆弧 MN，M 点和 N 点应在所求圆形磁场区域的边界上。

在通过 M、N 两点的不同的圆周中，最小的一个是以 MN 连线为直径的圆周。

\onepicture[w=6cm]{figs/31_an.jpg}

$r=\frac{1}{2}\mathrm{MN}=\frac{1}{2}\sqrt{R^{2}+R^{2}}=\frac{\sqrt{2}}{2}R=\frac{\sqrt{2}mv}{2qB}$ ②

所以本题所求的圆形磁场区域的最小半径为 $\frac{\sqrt{2}mv}{2qB}$。

所求磁场区域如图中实线圆所示。
}

\end{enumerate}

\end{document}
